\chapter{Fluencia y vida a alta temperatura}\label{ch:05}
\objetivos{Identificar los mecanismos de deformación dependiente del tiempo; utilizar leyes normalizadas de fluencia y distinguir extrapolación empírica de predicción validada.}
\section{Tiempo, temperatura y tensión}
La fluencia es deformación dependiente del tiempo bajo carga. En materiales metálicos suele adquirir relevancia a temperaturas homólogas elevadas, pero no existe una frontera universal que convierta cualquier comportamiento por encima de una fracción fija de la fusión en fluencia dominante. La temperatura homóloga $T/T_m$ se construye con temperaturas absolutas. En polímeros pueden observarse efectos viscoelásticos significativos a temperaturas mucho más bajas en términos absolutos.

En un ensayo a carga o tensión aproximadamente constante aparecen una deformación inicial y, con frecuencia, tres etapas: primaria, con velocidad decreciente; secundaria o casi estacionaria, con velocidad aproximadamente mínima; terciaria, con aceleración hasta rotura. No todos los ensayos muestran las tres de forma nítida. Bajo carga constante, el cambio de sección puede aumentar la tensión verdadera y contribuir a la aceleración.
\ilustra{fluencia-curva}{Evolución esquemática de deformación de fluencia a temperatura y carga mantenidas. Las regiones delimitan un comportamiento ilustrativo.}
\section{Mecanismos microestructurales}
El deslizamiento de dislocaciones permite deformación plástica; a temperatura suficiente, el ascenso facilita superar obstáculos mediante difusión de vacantes. La competencia entre endurecimiento y recuperación explica parte de la evolución hacia una tasa casi estacionaria. La difusión por el interior del grano y por sus límites puede producir deformación a tensiones relativamente bajas, con distinta dependencia de tamaño de grano.

El deslizamiento de límites necesita mecanismos de acomodación para evitar incompatibilidades. La cavitación puede iniciarse y crecer en límites o interfaces, reducir sección resistente y conducir a rotura. Una microestructura de grano fino, beneficiosa para ciertos comportamientos a temperatura ambiente, puede ofrecer más caminos de difusión o deslizamiento a alta temperatura. Esto ayuda a entender la solidificación direccional y los monocristales en álabes, sin afirmar que sean óptimos para cualquier pieza.

La fluencia coexiste con oxidación, fatiga térmica y cambios de precipitados. En una turbina, el arranque, la permanencia y la parada no son intercambiables con un ensayo a temperatura constante. Las interacciones dependen de secuencia y deben investigarse con datos adecuados.
\section{Ley de potencia y activación térmica}
Para un régimen de mecanismo dominante, una ley de Norton--Arrhenius puede escribirse
\begin{equation}\dot\varepsilon_s=A\sigma^n\exp\!\left(-\frac{Q}{RT}\right).\label{eq:norton}\end{equation}
Las unidades de $A$ dependen de la unidad de tensión. $Q$ y $R$ deben utilizar la misma base molar y unidades energéticas; $T$ es absoluta. El exponente $n$ y la energía aparente $Q$ proceden de ajustes en un dominio, no son constantes universales de todas las temperaturas.

Una forma más segura para comparar estados dentro del mismo régimen usa una referencia:
\begin{equation}\frac{\dot\varepsilon_2}{\dot\varepsilon_1}=\left(\frac{\sigma_2}{\sigma_1}\right)^n\exp\!\left[-\frac{Q}{R}\left(\frac1{T_2}-\frac1{T_1}\right)\right].\label{eq:nortonratio}\end{equation}
Las razones de tensiones son adimensionales y se reduce el riesgo de mezclar MPa y Pa en el coeficiente. Si $T_2>T_1$, el término entre paréntesis es negativo y la exponencial aumenta la tasa cuando $Q>0$: es una comprobación física del signo.
\caso{M05. Sensibilidad a un incremento térmico}{Un modelo ficticio de un único mecanismo tiene $Q=250\,\mathrm{kJ/mol}$ y mantiene la misma tensión al pasar de $900$ a $950\,\mathrm K$. La razón de tasas es $\exp[-250000/8.314462618\,(1/950-1/900)]\simeq5.803$. El resultado muestra una sensibilidad fuerte a la temperatura. No significa que cualquier superaleación aumente exactamente esa proporción: el valor de $Q$ y la permanencia del mecanismo son hipótesis del ejemplo.
\par\textbf{Desarrollo para comparar tu intento.} Al dividir la ley de Arrhenius a igual tensión, se cancela el prefactor y queda $\dot\varepsilon_2/\dot\varepsilon_1=\exp[-Q(1/T_2-1/T_1)/R]$. La diferencia de inversas es negativa al calentar, haciendo positivo el exponente.

\begin{enumerate}

\item Restar inversas de temperaturas absolutas.

\[\fitmath{\mathrm{inverseDifference}=\frac{1}{950}-\frac{1}{900}\approx -5.847953\times10^{-5}\,\mathrm{1/K}}\]

\item Convertir 250 kJ/mol a J/mol.

\[\fitmath{\mathrm{exponent}=\frac{-250000}{8.314462618}\cdot \left(\frac{1}{950}-\frac{1}{900}\right)\approx 1.758368\,}\]

\item Exponenciar para recuperar razón de tasas.

\[\fitmath{\mathrm{rateRatio}=\exp\left(\frac{-250000}{8.314462618}\cdot \left(\frac{1}{950}-\frac{1}{900}\right)\right)\approx 5.802958\,}\]

\end{enumerate}

}
\section{Estimación de deformación y relajación}
Si se aproxima constante la tasa secundaria, $\Delta\varepsilon\simeq\dot\varepsilon_s\Delta t$. Esta operación no calcula automáticamente deformación primaria, terciaria ni tiempo a rotura. Para una tasa $2\times10^{-8}\,\mathrm{s^{-1}}$ durante $100\,\mathrm h$, el incremento es $0.0072$, es decir, $0.72\,\%$; convertir horas a segundos es imprescindible.

En relajación de tensiones, la deformación total impuesta se mantiene aproximadamente constante mientras disminuye la parte elástica y aumenta la parte viscosa. Con elasticidad lineal, $0=\dot\sigma/E+\dot\varepsilon_c$, de donde $\dot\sigma=-E\dot\varepsilon_c$. Un ensayo de relajación no es equivalente a mantener tensión constante. Las juntas y uniones pueden perder precarga aun sin una rotura inmediata de la pieza.
\section{Larson--Miller y extrapolación}
Una correlación de tiempo a rotura emplea
\begin{equation}P=T\left(C_{LM}+\log_{10}t_r\right).\label{eq:larson}\end{equation}
La unidad de $t_r$ se fija en la calibración, habitualmente horas. Dentro del logaritmo se entiende el valor numérico respecto a esa unidad de referencia; no se cambia de horas a segundos conservando el mismo $C_{LM}$. El parámetro no elimina el efecto de la tensión: se construyen relaciones entre $P$ y tensión para material y estado determinados.

Igualar $P$ entre dos temperaturas a una misma tensión permite despejar un tiempo equivalente dentro de la correlación. Si cambian microestructura, mecanismo, oxidación o régimen de tensión, la extrapolación puede perder significado. Un ajuste de datos de ensayos cortos no prueba la vida durante decenas de miles de horas.
\section{Aplicación a aeromotores}
Los álabes calientes necesitan estabilidad microestructural y protección térmica; los discos, además, combinan gradientes, esfuerzos centrífugos y ciclos. El aumento de holguras por deformación puede deteriorar rendimiento antes de que exista rotura. La monitorización de temperatura del gas aporta información operativa, pero no mide directamente la temperatura local del metal ni el daño acumulado.

Para estudiar vida se conserva la historia térmica y mecánica, se definen zonas críticas y se utilizan datos representativos del proceso y del estado metalúrgico. Las inspecciones buscan indicadores compatibles con los mecanismos esperados. No se transforma una correlación académica en una extensión de vida de un componente sin el proceso técnico y documental aplicable.
\section{Monocristales y mecanismos de fluencia: aclaración de examen}
Un monocristal carece de límites de grano internos, pero no queda limitado a fluencia por difusión. También pueden existir mecanismos basados en dislocaciones dentro del grano, como deslizamiento y ascenso activados térmicamente. Es falsa la afirmación de que toda fluencia en un monocristal se produce por difusión, así como la idea de que la ausencia de límites de grano impide la fluencia.

Temperatura y tensión pueden modificar el mecanismo dominante; por ello una ley de Norton o una correlación de Larson--Miller solo se extrapola con justificación de su dominio. \emph{Tiempo para acumular deformación} no significa \emph{tiempo hasta rotura}: las etapas primaria y terciaria, el daño, la oxidación y la historia de cargas pueden alterar la predicción. En un álabe rotatorio, las tensiones centrífugas también requieren justificar distribución de masa, velocidad angular y sección resistente, en lugar de introducir una fuerza global arbitraria.


\section{Mecanismos, tamaño de grano y temperatura homóloga}
La temperatura homóloga $T/T_m$ compara temperaturas absolutas; sirve para orientar qué procesos activados pueden importar, sin fijar una frontera universal de fluencia. El deslizamiento de dislocaciones puede estar limitado por obstáculos; a mayor temperatura, el ascenso asistido por difusión permite superar ciertos obstáculos. Los mecanismos difusivos y el deslizamiento de bordes de grano pueden cobrar importancia con tensiones menores, granos finos y tiempos largos. Los mapas de mecanismos dependen de material y microestructura.

El afinamiento de grano puede mejorar resistencia a cedencia en ciertas condiciones de baja temperatura y, a la vez, aumentar la contribución de mecanismos ligados a bordes a alta temperatura. No hay contradicción: se trata de mecanismos y condiciones diferentes. Esta competencia ayuda a entender el interés de granos orientados o monocristales en ciertos álabes, sin deducir que eliminen todos los procesos de deformación y daño.
\section{Cómo identificar una ley de Norton--Arrhenius}
En un dominio con $\dot\varepsilon=A\sigma^n\exp[-Q/(RT)]$, comparar ensayos a la misma temperatura permite estimar $n$ de la pendiente de $\ln\dot\varepsilon$ frente a $\ln\sigma$. A tensión fija, una gráfica de $\ln\dot\varepsilon$ frente a $1/T$ tiene pendiente $-Q/R$, si $A$ y el mecanismo se mantienen. Un cambio de pendiente puede señalar que una sola ley no describe todo el conjunto.

En una formulación normalizada se escriben cocientes respecto a un estado de referencia; así se reduce el riesgo de confundir unidades del prefactor. Aun así, una buena coincidencia del ajuste no prueba la estabilidad futura de precipitados, oxidación o ausencia de cavitación. El intervalo de calibración debe acompañar a los parámetros.
\Needspace{.48\textheight}
\section{Relajación bajo deformación impuesta}
Bajo deformación total constante y temperatura fija, $\varepsilon_{tot}=\sigma/E+\varepsilon_c$ conduce, si $E$ es constante, a $\dot\sigma=-E\dot\varepsilon_c$. Para Norton a esa temperatura,
\begin{equation}\frac{d\sigma}{dt}=-EA_T\sigma^n,\qquad \sigma(t)^{1-n}=\sigma_0^{1-n}+(n-1)EA_Tt\quad(n\ne1).\end{equation}
Para $n=1$ resulta una disminución exponencial. La tensión se relaja al transferirse deformación elástica a deformación de fluencia. Es distinto de un ensayo a carga constante, donde la tensión no se reduce por esa misma condición de contorno. Con cambios de sección o temperatura hacen falta términos y modelos adicionales.
\repaso{¿Por qué una tasa secundaria baja no demuestra una vida infinita?}{Porque la tasa describe un régimen de deformación, no todo el daño. Pueden evolucionar microestructura, cavidades, oxidación y distribución de tensiones, y puede aparecer fluencia terciaria antes de rotura.}

\resumen{La tasa de fluencia y el tiempo a rotura son magnitudes distintas. La temperatura de las leyes de activación es absoluta. Las constantes de una ley dimensional dependen de las unidades. Una extrapolación solo es defendible dentro de un régimen y con evidencia de su continuidad.}
\fuentes{Roylance, \textit{Dislocation Basis of Yield and Creep} \cite{MIT-CREEP}; materiales docentes de fluencia, pp. 1--27; investigación de superaleaciones de NASA \cite{NASA-SX}.}
